AWS PrivateLink endpoint có một địa chỉ IP cố định trong từng AZ nơi
chúng được triển khai, trong suốt thời gian tồn tại của endpoint (cho
đến khi endpoint bị xóa). Các địa chỉ IP này được gắn vào Elastic
network interface (ENI). AWS khuyến nghị sử dụng DNS để resolve địa chỉ
IP cho endpoint để các ứng dụng downstream sử dụng địa chỉ IP mới nhất
khi ENIs được thêm vào AZ mới hoặc bị xóa theo thời gian.

Trong phần này, bạn sẽ tạo một quy tắc chuyển tiếp (forwarding rule) để
gửi các yêu cầu resolve DNS từ môi trường truyền thống (mô phỏng) đến
Private Hosted Zone trên Route 53. Phần này tận dụng cơ sở hạ tầng do
CloudFormation triển khai trong phần Chuẩn bị môi trường.

\vspace{0.5em}\noindent\textbf{Tạo DNS Alias Records cho Interface
endpoint}\par

\begin{enumerate}
\def\labelenumi{\arabic{enumi}.}
\tightlist
\item
  Click link để đi đến
  \href{https://us-east-1.console.aws.amazon.com/route53/v2/hostedzones?region=us-east-1\#}{Route
  53 management console} (Hosted Zones section). Mẫu CloudFormation mà
  bạn triển khai trong phần Chuẩn bị môi trường đã tạo Private Hosted
  Zone này. Nhấp vào tên của Private Hosted Zone,
  s3.us-east-1.amazonaws.com:
\end{enumerate}

\begin{figure}
\centering
\pandocbounded{\includegraphics[keepaspectratio,alt={hosted zone}]{5-Workshop/5.4-S3-onprem/hosted-zone.png}}
\caption{hosted zone}
\end{figure}

\begin{enumerate}
\def\labelenumi{\arabic{enumi}.}
\setcounter{enumi}{1}
\tightlist
\item
  Tạo 1 record mới trong Private Hosted Zone:
\end{enumerate}

\begin{figure}
\centering
\pandocbounded{\includegraphics[keepaspectratio,alt={Create record}]{5-Workshop/5.4-S3-onprem/create-record1.png}}
\caption{Create record}
\end{figure}

\begin{itemize}
\tightlist
\item
  Giữ nguyên Record name và record type
\item
  Alias Button: click để enable
\item
  Route traffic to: Alias to VPC Endpoint
\item
  Region: US East (N. Virginia) {[}us-east-1{]}
\item
  Chọn endpoint: Paste tên (Regional VPC Endpoint DNS) bạn đã lưu lại ở
  phần 4.3
\end{itemize}

\begin{figure}
\centering
\pandocbounded{\includegraphics[keepaspectratio,alt={record1}]{5-Workshop/5.4-S3-onprem/record1.png}}
\caption{record1}
\end{figure}

\begin{enumerate}
\def\labelenumi{\arabic{enumi}.}
\setcounter{enumi}{2}
\tightlist
\item
  Click Add another record, và add 1 cái record thứ 2 sử dụng những
  thông số sau:
\end{enumerate}

\begin{itemize}
\tightlist
\item
  Record name: *.
\item
  Record type: giữ giá trị default (type A)
\item
  Alias Button: Click để enable
\item
  Route traffic to: Alias to VPC Endpoint
\item
  Region: US East (N. Virginia) {[}us-east-1{]}
\item
  Chọn endpoint: Paste tên (Regional VPC Endpoint DNS) bạn đã lưu lại ở
  phần 4.3
\item
  Click \textbf{Create records}
\end{itemize}

\begin{figure}
\centering
\pandocbounded{\includegraphics[keepaspectratio,alt={record 2}]{5-Workshop/5.4-S3-onprem/record2.png}}
\caption{record 2}
\end{figure}

Record mới xuất hiện trên giao diện Route 53.

\begin{figure}
\centering
\pandocbounded{\includegraphics[keepaspectratio,alt={result}]{5-Workshop/5.4-S3-onprem/result.png}}
\caption{result}
\end{figure}

\vspace{0.5em}\noindent\textbf{Tạo một Resolver Forwarding
Rule}\par

\textbf{Route 53 Resolver Forwarding Rules} cho phép bạn chuyển tiếp các
DNS queries từ VPC của bạn đến các nguồn khác để resolve name. Bên ngoài
môi trường workshop, bạn có thể sử dụng tính năng này để chuyển tiếp các
DNS queries từ VPC của bạn đến các máy chủ DNS chạy trên on-premises.
Trong phần này, bạn sẽ mô phỏng một on-premises conditional forwarder
bằng cách tạo một forwarding rule để chuyển tiếp các DNS queries for
Amazon S3 đến một Private Hosted Zone chạy trong ``VPC Cloud'' để
resolve PrivateLink interface endpoint regional DNS name.

\begin{enumerate}
\def\labelenumi{\arabic{enumi}.}
\item
  Từ giao diện \textbf{Route 53}, chọn \textbf{Inbound endpoints} trên
  thanh bên trái
\item
  Trong giao diện \textbf{Inbound endpoint}, Chọn ID của Inbound
  endpoint.
\end{enumerate}

\begin{figure}
\centering
\pandocbounded{\includegraphics[keepaspectratio,alt={Inbound endpoint}]{5-Workshop/5.4-S3-onprem/route53-1.png}}
\caption{Inbound endpoint}
\end{figure}

\begin{enumerate}
\def\labelenumi{\arabic{enumi}.}
\setcounter{enumi}{2}
\tightlist
\item
  Sao chép 2 địa chỉ IP trong danh sách vào trình chỉnh sửa.
\end{enumerate}

\begin{figure}
\centering
\pandocbounded{\includegraphics[keepaspectratio,alt={Ip addresses}]{5-Workshop/5.4-S3-onprem/route53-2.png}}
\caption{Ip addresses}
\end{figure}

\begin{enumerate}
\def\labelenumi{\arabic{enumi}.}
\setcounter{enumi}{3}
\tightlist
\item
  Từ giao diện Route 53, chọn \textbf{Resolver} \textgreater{}
  \textbf{Rules} và chọn \textbf{Create rule}
\end{enumerate}

\begin{figure}
\centering
\pandocbounded{\includegraphics[keepaspectratio,alt={Ip addresses}]{5-Workshop/5.4-S3-onprem/route53-3.png}}
\caption{Ip addresses}
\end{figure}

\begin{enumerate}
\def\labelenumi{\arabic{enumi}.}
\setcounter{enumi}{4}
\tightlist
\item
  Trong giao diện \textbf{Create rule}
\end{enumerate}

\begin{itemize}
\tightlist
\item
  Name: myS3Rule
\item
  Rule type: Forward
\item
  Domain name: s3.us-east-1.amazonaws.com
\end{itemize}

\begin{figure}
\centering
\pandocbounded{\includegraphics[keepaspectratio,alt={create rule}]{5-Workshop/5.4-S3-onprem/route53-4.png}}
\caption{create rule}
\end{figure}

\begin{itemize}
\tightlist
\item
  VPC: VPC On-prem
\item
  Outbound endpoint: VPCOnpremOutboundEndpoint
\end{itemize}

\begin{figure}
\centering
\pandocbounded{\includegraphics[keepaspectratio,alt={create rule}]{5-Workshop/5.4-S3-onprem/route53-5.png}}
\caption{create rule}
\end{figure}

\begin{itemize}
\tightlist
\item
  Target IP Addresses: điền cả hai IP bạn đã lưu trữ trên trình soạn
  thảo (inbound endpoint addresses) và sau đó chọn \textbf{Submit}
\end{itemize}

\begin{figure}
\centering
\pandocbounded{\includegraphics[keepaspectratio,alt={create rule}]{5-Workshop/5.4-S3-onprem/route53-6.png}}
\caption{create rule}
\end{figure}

Bạn đã tạo thành công resolver forwarding rule.

\begin{figure}
\centering
\pandocbounded{\includegraphics[keepaspectratio,alt={create rule}]{5-Workshop/5.4-S3-onprem/route53-7.png}}
\caption{create rule}
\end{figure}

\vspace{0.5em}\noindent\textbf{Kiểm tra on-premises DNS mô
phỏng.}\par

\begin{enumerate}
\def\labelenumi{\arabic{enumi}.}
\tightlist
\item
  Kết nối đến \textbf{Test-Interface-Endpoint EC2 instance} với
  \textbf{Session Manager}
\end{enumerate}

\begin{figure}
\centering
\pandocbounded{\includegraphics[keepaspectratio,alt={create rule}]{5-Workshop/5.4-S3-onprem/test1.png}}
\caption{create rule}
\end{figure}

\begin{enumerate}
\def\labelenumi{\arabic{enumi}.}
\setcounter{enumi}{1}
\tightlist
\item
  Kiểm tra DNS resolution. Lệnh dig sẽ trả về địa chỉ IP được gán cho
  VPC endpoint interface đang chạy trên VPC (địa chỉ IP của bạn sẽ
  khác):
\end{enumerate}

\begin{verbatim}
dig +short s3.us-east-1.amazonaws.com 
\end{verbatim}

Các địa chỉ IP được trả về là các địa chỉ IP VPC enpoint, KHÔNG phải là
các địa chỉ IP Resolver mà bạn đã dán từ trình chỉnh sửa văn bản của
mình. Các địa chỉ IP của Resolver endpoint và VPC endpoin trông giống
nhau vì chúng đều từ khối CIDR VPC Cloud.

\begin{figure}
\centering
\pandocbounded{\includegraphics[keepaspectratio,alt={create rule}]{5-Workshop/5.4-S3-onprem/dig.png}}
\caption{create rule}
\end{figure}

\begin{enumerate}
\def\labelenumi{\arabic{enumi}.}
\setcounter{enumi}{2}
\tightlist
\item
  Truy cập vào menu VPC (phần Endpoints), chọn S3 interface endpoint.
  Nhấp vào tab Subnets và xác nhận rằng các địa chỉ IP được trả về bởi
  lệnh Dig khớp với VPC endpoint:
\end{enumerate}

\begin{figure}
\centering
\pandocbounded{\includegraphics[keepaspectratio,alt={create rule}]{5-Workshop/5.4-S3-onprem/subnet.png}}
\caption{create rule}
\end{figure}

\begin{enumerate}
\def\labelenumi{\arabic{enumi}.}
\setcounter{enumi}{3}
\tightlist
\item
  Hãy quay lại shell của bạn và sử dụng AWS CLI để kiểm tra danh sách
  các bucket S3 của bạn:
\end{enumerate}

\begin{verbatim}
aws s3 ls --endpoint-url https://s3.us-east-1.amazonaws.com
\end{verbatim}

\begin{figure}
\centering
\pandocbounded{\includegraphics[keepaspectratio,alt={create rule}]{5-Workshop/5.4-S3-onprem/endpoint.png}}
\caption{create rule}
\end{figure}

\begin{enumerate}
\def\labelenumi{\arabic{enumi}.}
\setcounter{enumi}{4}
\tightlist
\item
  Kết thúc phiên làm việc của Session Manager của bạn:
\end{enumerate}

\begin{figure}
\centering
\pandocbounded{\includegraphics[keepaspectratio,alt={create rule}]{5-Workshop/5.4-S3-onprem/terminal.png}}
\caption{create rule}
\end{figure}

Trong phần này, bạn đã tạo một \textbf{Interface Endpoint} cho Amazon
S3. Điểm cuối này có thể được truy cập từ on-premises thông qua
Site-to-Site VPN hoặc AWS Direct Connect. Các điểm cuối Route 53
Resolver outbound giả lập chuyển tiếp các yêu cầu DNS từ on-premises đến
một Private Hosted Zone đang chạy trên đám mây. Các điểm cuối Route 53
inbound nhận yêu cầu giải quyết và trả về một phản hồi chứa địa chỉ IP
của \textbf{Interface Endpoint} VPC. Sử dụng DNS để giải quyết các địa
chỉ IP của điểm cuối cung cấp tính sẵn sàng cao trong trường hợp một
Availability Zone gặp sự cố.
